\begin{center}
\begin{minipage}{0.98\linewidth}
\begin{tcolorbox}[colback=white,colframe=CStatement!25,boxrule=0.4pt,
  arc=1mm,left=3mm,right=3mm,top=3mm,bottom=3mm]
\centering
\textbf{对称圆：圆心移到对称点，半径保持不变}\par\medskip
\begin{tikzpicture}[font=\small,>=Stealth]

\begin{scope}[x=.78cm,y=.78cm]
\coordinate (C) at (1,2);\coordinate (D) at (.2,.4);\coordinate (H) at (.6,1.2);
\draw[CExample,thick] (C) circle (3);
\draw[CTrap,thick,dashed] (D) circle (3);
\draw[CStatement,thick] (-3,3)--(5,-1) node[right] {$l$};
\draw[CProof,thick] (C)--(D);
\draw[CExample,thick] (C)--(1,5);
\node[right,text=CExample] at (1,4.2) {$3$};
\draw[CTrap,thick] (D)--(.2,-2.6);
\node[left,text=CTrap] at (.2,-1.8) {$3$};
\fill[CExample] (C) circle (2.5pt);\fill[CTrap] (D) circle (2.5pt);
\node[right] at (C) {$C$};\node[left] at (D) {$C'$};
\draw[CProof] (.68,1.36)--(.84,1.28)--(.76,1.12);
\end{scope}
\end{tikzpicture}\par\medskip
实线圆心 $C(1,2)$，虚线圆心 $C'(1/5,2/5)$，两圆半径均为 $3$；对称轴为 $l:x+2y-3=0$。
\end{tcolorbox}
\end{minipage}
\end{center}
